\section{Introduction}
%----------------------------------------------------------------------
The behaviour of the strong nuclear field is conventionally described within
Quantum Chromodynamics. Here we adopt a different point of departure: a single
geometric assumption---that the three generating circles of the Reuleaux triangle
are retained as an active part of the construction. From this assumption a
continuous chain follows: the Force Angle Law, the Angular Shortfall, the
four-sphere membrane body, the layered particle, residual stress, helical stacking,
optical organisation of the residual field, and cosmic cyclicity.

The present version substantially extends earlier drafts. In addition to the
geometric skeleton we now include: (i)~the residual-stress dipole and its
combinatorial necessity; (ii)~the pressure length as a pure geometric distortion
scale; (iii)~elastic moduli ratios fixed by the tetrahedron; (iv)~a five-mode
optical taxonomy controlled by the shortfall; (v)~dual dissipation channels;
(vi)~absolute scale closure with $\hbar c$ and a hadronic choice of $s$; and
(vii)~refined aggregation and dissolution pathways at $N_c^*=3$.
